\chapter{序列建模模型}
\begin{quote}\small
\textit{本章摘要。}序列模型根据有顺序的历史观测形成表示，再由任务头预测未来数值或事件。异常检测还需要将已到达的观测与正常模式比较。本章先固定这几类任务的信息边界，再围绕“怎样利用历史”讲解代表模型：RNN 与 LSTM 逐步更新记忆，时间卷积网络组合固定范围的历史，Transformer 按内容读取历史，Mamba 用输入依赖的状态更新处理长序列。最后将这些表示接到多步预测、概率预测和异常评分上，说明输出如何训练、如何评价以及何时能够触发告警。
\end{quote}
设每分钟记录一次温度、转速和振动。过去一小时的平均温度相同，可能对应稳定运行，也可能对应持续升温后骤降。平均值没有保留这些变化的顺序，而未来温度和异常判断可能都依赖这一顺序。序列模型的共同工作，就是把有顺序的观测变成对当前任务有用的表示。

本章以设备监测为例，先用未来温度预测走通输入、表示和输出，再说明同一类表示如何用于异常检测与未来故障预测。模型部分围绕历史信息的处理方式展开。逐步压缩历史带来记忆与训练问题，固定时间范围的卷积提供另一种计算方式，注意力允许重新读取远处的记录，选择性状态更新则控制长历史的计算和存储成本。它们是不同约束下的选择，不构成必然逐级升级的关系。

\section{先固定任务与可用的历史}
\subsection{预测未来与识别当前异常}
设$x_t\in\R^p$为时刻$t$已经取得的$p$个观测，长度为$T$的输入窗口为$X_t=(x_{t-T+1},\ldots,x_t)$。温度预测要根据这个窗口估计未来数值，例如未来$H$分钟的温度$(y_{t+1},\ldots,y_{t+H})$。故障预测则估计未来指定时间段内发生故障的概率，需要按这一时间段定义零一事件标签。剩余寿命预测估计距失效还剩多久，又是另一种目标。它们可以使用同一类序列编码器，但输出含义与标签不同。

异常检测判断已经到达的观测是否偏离正常模式。例如，在$t-1$时预测温度$\hat y_{t\mid t-1}$，等$y_t$到达后才计算误差。这个误差可以提示当前异常，但不直接等于未来故障概率。若要用它提前预警，必须另外验证异常与后续故障之间的关系。

\subsection{共享表示与不同的输出头}
把序列变成向量或向量序列的网络主体称为\textbf{序列编码器}，把表示转成指定预测的部分称为\textbf{任务头}。可将窗口预测写为
\begin{equation}
 r_t=f_\theta(X_t),\qquad \hat y_{t+1:t+H}=g_\phi(r_t).
 \label{eq:sequence-encoder-head}
\end{equation}
$r_t$是中间表示，$\theta$与$\phi$分别为编码器和任务头参数。数值预测头可以输出未来温度，事件预测头可以输出故障概率。异常评分还可以使用编码器产生的正常模式预测，再与实际观测比较。共享的是表示方法，训练目标不同仍可能使编码器学到不同参数，不能假定一个训练好的表示适合所有任务。

图\ref{fig:sequence-tasks}区分编码、输出和事后比较三种操作。图中的实线计算在时刻$t-1$完成，虚线在新观测到达后才形成残差。未来故障头须用事件标签训练，不能通过给残差改名获得。
\begin{figure}[htbp]
\centering
\begin{tikzpicture}[node distance=0.8cm and 0.8cm]
 \node[idi input,minimum width=2.0cm] (x) {已知窗口\\$X_{t-1}$};
 \node[idi state,right=of x,minimum width=2.1cm] (r) {编码器表示\\$r_{t-1}$};
 \node[idi output,right=of r,minimum width=2.5cm] (p) {温度预测\\$\hat y_{t\mid t-1}$};
 \node[idi output,below=of p,minimum width=2.5cm] (risk) {未来事件概率\\$p_{t-1}$};
 \node[idi input,above=of p,minimum width=2.5cm] (obs) {到达的温度\\$y_t$};
 \node[idi process,right=of p,minimum width=2.2cm] (score) {比较后评分\\$|y_t-\hat y_{t\mid t-1}|$};
 \draw[idi arrow] (x)--(r);\draw[idi arrow] (r)--(p);
 \draw[idi arrow] (r.south)|-(risk.west);
 \draw[idi arrow,dashed] (p)--(score);\draw[idi arrow,dashed] (obs.east)-|(score.north);
\end{tikzpicture}
\caption{相同类别的历史编码方法可以服务于不同任务。实线是截至$t-1$的预测计算，虚线是$t$时刻观测到达后的异常评分。各任务可分别训练，也可在有相应标签时联合训练。}
\label{fig:sequence-tasks}
\end{figure}

例如，四分钟的三个传感器组成$4\times3$输入，八台设备构成一批时为$8\times4\times3$。预测未来两分钟温度时，每个样本的标签是长度为二的向量。训练用已经记录下来的未来标签计算损失，预测时只能读取当前已取得的窗口。不同设备之间不传递状态。若未来转速来自已知排程，可以把计划值作为外生输入，否则未来真实转速不能提前进入模型。

\section{RNN 与 LSTM：把历史逐步写入记忆}
\subsection{RNN 的状态怎样形成预测}
当数据持续到达时，可以每读一个观测就更新一个固定宽度的状态，无需每次重新计算完整窗口。循环神经网络（Recurrent Neural Network，RNN）将这一规则写为
\begin{equation}
 h_t=\tanh(W_hh_{t-1}+W_xx_t+b),\qquad
 \hat y_{t+1}=w^{\mathsf T}h_t+b_y.
 \label{eq:sequence-rnn-predict}
\end{equation}
设$h_t\in\R^m$，则$W_x\in\R^{m\times p}$、$W_h\in\R^{m\times m}$、$b,w\in\R^m$，$b_y$为标量。$W_x$处理当前输入，$W_h$处理旧状态，同一组参数在各时刻重复使用。$\tanh$逐分量施加非线性。状态是内部计算的一组数，不是另一个传感器，也不是自动获得的状态概率分布。

图\ref{fig:sequence-recurrent-state}中，纵向箭头读入当前观测，横向箭头传递过去的状态。若$m=2$，每一步都把三个新观测与两个旧状态数变成两个新状态数。初始状态可以为零，也可以是学习到的共同初值。预测用平方误差等目标训练，梯度通过任务头和历史状态调整共享参数。

\begin{figure}[htbp]
\centering
\begin{tikzpicture}[node distance=0.55cm and 0.6cm]
 \node[idi input, minimum width=1.45cm] (x1) {$x_{t-2}$};
 \node[idi input, right=of x1, minimum width=1.45cm] (x2) {$x_{t-1}$};
 \node[idi input, right=of x2, minimum width=1.45cm] (x3) {$x_t$};
 \node[idi state, above=0.75cm of x1, minimum width=1.45cm] (h1) {$h_{t-2}$};
 \node[idi state, above=0.75cm of x2, minimum width=1.45cm] (h2) {$h_{t-1}$};
 \node[idi state, above=0.75cm of x3, minimum width=1.45cm] (h3) {$h_t$};
 \draw[idi arrow] (x1)--(h1); \draw[idi arrow] (x2)--(h2); \draw[idi arrow] (x3)--(h3);
 \draw[idi arrow] (h1)--(h2); \draw[idi arrow] (h2)--(h3);
 \node[above=0.25cm of h2,font=\scriptsize] {状态沿时间向前传递并压缩历史};
 \node[idi output, right=0.7cm of h3, minimum width=1.45cm] (y) {$\hat y_{t+1}$};
 \draw[idi arrow] (h3)--(y);
\end{tikzpicture}
\caption{递归序列模型把历史压缩为状态；状态维度和更新机制决定哪些信息能够持续保留。}
\label{fig:sequence-recurrent-state}
\end{figure}
状态宽度不随历史增长，但过去的信息必须经过连续更新才能影响当前输出。单次冲击既可能被后续状态覆盖，也可能在训练时难以获得有效梯度。下面先看门控如何改变更新，再用梯度路径说明它解决了什么困难。

\subsection{LSTM：分开控制保留、写入与读出}
长短期记忆网络（Long Short-Term Memory，LSTM）在状态更新中增加门控与单独的细胞状态\cite{hochreiter1997lstm}。它让网络分别控制旧内容的保留和新内容的写入：
\begin{align}
 i_t&=\sigma(W_ix_t+U_ih_{t-1}),& f_t&=\sigma(W_fx_t+U_fh_{t-1}),\label{eq:sequence-lstm-input-forget}\\
 o_t&=\sigma(W_ox_t+U_oh_{t-1}),& \tilde c_t&=\tanh(W_cx_t+U_ch_{t-1}),\label{eq:sequence-lstm-output-candidate}\\
 c_t&=f_t\odot c_{t-1}+i_t\odot\tilde c_t,& h_t&=o_t\odot\tanh(c_t).
\end{align}
这里$\sigma(a)=1/(1+e^{-a})$把数压到零与一之间，$\tanh$把数压到$(-1,1)$，$\odot$表示对应分量分别相乘，不是矩阵乘法。输入门$i_t$决定写多少候选内容，遗忘门$f_t$决定留多少旧内容，输出门$o_t$决定公开多少记忆；细胞状态$c_t$是存储通道，$h_t$是供下一步和任务头使用的表示。

可先手算一个分量：若$c_{t-1}=2,f_t=0.9,i_t=0.2,\tilde c_t=0.5,o_t=1/2$，则$c_t=0.9\times2+0.2\times0.5=1.9$，$h_t=\tanh(1.9)/2\approx0.478$。候选值受限不代表累积记忆也限于$[-1,1]$。这些门值是网络根据输入算出的，不是人为判定“重要”后再开门。细胞状态在保留旧值与加入新值之间提供了明确的加法路径。门控循环单元（Gated Recurrent Unit，GRU）进一步合并部分记忆通道，后面将与 LSTM 比较其更新方式。

门控的必要性可以从训练梯度理解。
对简单RNN，隐藏状态对早期状态的导数包含
\begin{equation}
 \frac{\partial h_T}{\partial h_t^{\mathsf T}}=J_TJ_{T-1}\cdots J_{t+1},\qquad
 J_s=\diag\left(1-\tanh^2(a_s)\right)W_h.
\end{equation}
这里$a_s=W_hh_{s-1}+W_xx_s+b$，Jacobian（雅可比矩阵）的第$(i,j)$项为输出第$i$分量对输入第$j$分量的偏导；矩阵因子按式中次序排列不可交换，作用于早期扰动列向量时则由最右侧因子先起作用；反向传播列梯度时使用这一 Jacobian 的转置。
若乘积中的最大奇异值长期小于一，梯度趋向零；若大于一，梯度可能爆炸。梯度裁剪在$g\ne0$时把$g$替换为$g\min(1,\tau/\|g\|)$，其中阈值$\tau>0$；$g=0$时直接保留零，避免除零。它可以控制爆炸但不能修复长期消失。LSTM的细胞状态通过门$f_t$保留一条近似加法路径，使信息保存时间由门值控制。

非线性导数$1-\tanh^2(a_s)$在饱和区域接近零，因此仅控制循环权重的尺度仍不足以保证长期梯度。对于长流，可以每$K$步截断计算图以控制训练内存，前向状态仍可携带更早的信息，但早期输入不能通过被截断的路径直接贡献梯度。窗口长度与截断长度因而是不同的实验变量。实现截断时，应保留状态数值但切断它与上一计算图的梯度连接；切断梯度不等于把状态清零。若下一批已换设备或不再时间连续，则应按样本边界重置，不能沿用上一批状态。

\paragraph{从单步导数到跨时间反向传播。}
设$x_t\in\R^p$、$h_t\in\R^m$，于是$W_x\in\R^{m\times p}$、$W_h\in\R^{m\times m}$、$b\in\R^m$。令总损失为$\mathcal L=\sum_{t=1}^T\ell_t(h_t)$，用列向量$g_t=\nabla_{h_t}\mathcal L$表示包含所有后续损失的梯度。链式法则逐步给出
\begin{align}
 g_T&=\nabla_{h_T}\ell_T,\\
 g_t&=\nabla_{h_t}\ell_t+J_{t+1}^{\mathsf T}g_{t+1},\\
 g_t&=\sum_{s=t}^T(J_sJ_{s-1}\cdots J_{t+1})^{\mathsf T}\nabla_{h_s}\ell_s,
\end{align}
其中$s=t$时空乘积为$I$。再令$\delta_t=\diag(1-h_t\odot h_t)g_t$，共享权重的梯度不是某一步的梯度，而是所有出现位置的累加：
\begin{equation}
 \nabla_{W_h}\mathcal L=\sum_{t=1}^T\delta_t h_{t-1}^{\mathsf T},\qquad
 \nabla_{W_x}\mathcal L=\sum_{t=1}^T\delta_t x_t^{\mathsf T}.
\end{equation}
由矩阵范数的次乘性，$\|J_s\cdots J_{t+1}\|_2\leq\prod_{r=t+1}^s\|J_r\|_2$。若每项至多为$\rho<1$，远期损失传回的上界按$\rho^{s-t}$衰减；若上界大于一，仅意味着可能放大，不能据此断言实际梯度必然爆炸，因为奇异方向会随时间改变。

在LSTM中，$x_t\in\R^p$，各门、$c_t,h_t\in\R^m$，所以各$W$为$m\times p$、各$U$为$m\times m$；式\eqref{eq:sequence-lstm-input-forget}和式\eqref{eq:sequence-lstm-output-candidate}省略了可训练偏置。固定门值时，沿细胞状态直接路径有
\begin{equation}
 \left.\frac{\partial c_T}{\partial c_t^{\mathsf T}}\right|_{\text{直接路径}}
 =\diag\left(\prod_{s=t+1}^T f_s\right).
\end{equation}
乘积按元素计算。若某维遗忘门近似恒为$f\in(0,1)$，经过$k$步保留$f^k$，其衰减到初值$e^{-1}$所需步数为$-1/\log f$。例如$f=0.9$时这一时间尺度约为$9.49$步，$f=0.99$时约为$99.50$步；它是衰减到约$36.8\%$的时间，不是到零的时间。这解释了门值与记忆时间尺度的关系，但不是完整Jacobian：门还通过$h_{t-1}$依赖早期状态，其他路径仍会放大或衰减梯度。门控改善可训练性，不保证无限记忆。

\subsection{GRU：用较少门控更新状态}

GRU使用更新门和重置门
\begin{align}
 z_t&=\sigma(W_zx_t+U_zh_{t-1}),\\
 r_t&=\sigma(W_rx_t+U_rh_{t-1}),\\
 \tilde h_t&=\tanh(W_hx_t+U_h(r_t\odot h_{t-1})),\\
 h_t&=(1-z_t)\odot h_{t-1}+z_t\odot\tilde h_t.
\end{align}
更新门接近零时保留旧状态，接近一时写入候选状态。它比LSTM参数更少，却牺牲了对细粒度记忆通道的独立控制。

\section{时间卷积网络：在固定范围内组合历史}
\subsection{因果卷积怎样保留顺序}
递归状态按时间逐步计算。若每次输入都是一段完整窗口，也可以同时计算窗口中各位置的局部特征。时间卷积网络（Temporal Convolutional Network，TCN）沿时间轴使用因果卷积，并通过多层空洞卷积扩大可见范围。因果卷积在位置$t$只读取当前和过去，单层可写为
\begin{equation}
 z_t=\rho\left(b+\sum_{j=0}^{k-1}W_jx_{t-jd}\right).
 \label{eq:sequence-tcn}
\end{equation}
$k$为卷积核宽度，$d$为正整数空洞率，$W_j$把输入通道映射到输出通道，$\rho$为逐分量激活函数。边界以零补齐并只在左侧补值。$d=1$读取连续位置，$d=2$每隔一个位置读取一次。各相对位置有各自权重，所以交换两个时间点通常会改变结果。

以标量输入为例，令$k=2,d=2,W_0=1,W_1=-1,b=0$，暂不使用非线性，则$z_t=x_t-x_{t-2}$。温度序列$(20,21,22)$在第三个位置给出$2$，交换为$(22,21,20)$后给出$-2$。两段的均值相同，卷积提取的变化方向不同。多通道卷积可以同时组合温度、转速与振动，实际权重由预测损失学习。

\subsection{空洞卷积怎样扩大感受野}
设每层步幅为一、核宽为$k$，第$\ell$层空洞率为$d_\ell$，则$L$层输出能依赖的最早位置与当前位置之间跨度为
\begin{equation}
 R=1+(k-1)\sum_{\ell=1}^L d_\ell.
\end{equation}
$R$是感受野跨度，不自动保证跨度内每个位置都被覆盖。取$k=2$、$d_\ell=2^{\ell-1}$时，路径可组合这些二进制偏移，覆盖最近$2^L$个位置。三层空洞率$1,2,4$便可覆盖八步历史。若每个残差块有两层同空洞率卷积，计算感受野时须将两层都计入。

图\ref{fig:sequence-tcn-receptive}把三层的局部读取展开：最终节点先取相隔四步的两个表示，每个表示再取相隔两步的两个表示，底层则读取相邻输入。
\begin{figure}[htbp]
\centering
\begin{tikzpicture}[x=1.05cm,y=0.8cm,>=Latex,every node/.style={font=\scriptsize}]
 \foreach \i in {0,...,7} {\node[draw,circle,inner sep=2pt] (x\i) at (\i,0) {$x_{t-7+\i}$};}
 \foreach \i in {0,...,3} {\node[draw,circle,inner sep=3pt] (a\i) at (2*\i+0.5,1) {};}
 \foreach \i/\j/\k in {0/0/1,1/2/3,2/4/5,3/6/7} {\draw[->] (x\j)--(a\i);\draw[->] (x\k)--(a\i);}
 \node[draw,circle,inner sep=3pt] (b0) at (1.5,2) {};
 \node[draw,circle,inner sep=3pt] (b1) at (5.5,2) {};
 \draw[->] (a0)--(b0);\draw[->] (a1)--(b0);
 \draw[->] (a2)--(b1);\draw[->] (a3)--(b1);
 \node[draw,circle,inner sep=3pt] (c) at (3.5,3) {$z_t$};
 \draw[->] (b0)--(c);\draw[->] (b1)--(c);
 \node[right] at (7.4,1) {$d=1$};\node[right] at (7.4,2) {$d=2$};\node[right] at (7.4,3) {$d=4$};
\end{tikzpicture}
\caption{核宽为二的三层因果空洞卷积中，当前位置输出读取最近八步的路径。箭头表示特征参与下一层计算，只展开与$z_t$相关的节点。}
\label{fig:sequence-tcn-receptive}
\end{figure}

TCN 常用残差块，将卷积变换结果与输入相加，通道宽度不同时先投影。增加层数或空洞率可以扩大历史范围，但超出感受野的信息仍无法影响当前输出。训练时各位置可并行计算，流式运行也可以缓存各层所需的过去特征。它适合有明确局部时间尺度的预测，不能仅凭扩大跨度保证保留短冲击或远期周期。下面讨论另一种做法：保留历史表示，在预测需要时重新读取。

\section{Transformer：按内容重新读取历史}
\subsection{查询、键和值怎样完成一次读取}
图\ref{fig:sequence-attention-query}区分历史键用于匹配、历史值用于输出的两条路径；当前查询通过归一化权重决定读取哪些历史内容。
\begin{figure}[htbp]
\centering
\begin{tikzpicture}[node distance=0.8cm and 0.7cm]
 \node[idi input, minimum width=1.8cm] (q) {当前查询\\$q_i$};
 \node[idi state, right=of q, minimum width=2.0cm] (score) {与历史键\\计算相似度};
 \node[idi process, right=of score, minimum width=2.0cm] (weight) {softmax 权重\\$a_{ij}$};
 \node[idi output, right=of weight, minimum width=2.0cm] (out) {加权值\\$o_i$};
 \draw[idi arrow] (q)--(score); \draw[idi arrow] (score)--(weight); \draw[idi arrow] (weight)--(out);
 \node[idi input, below=0.75cm of score, minimum width=2.0cm] (memory) {历史键和值\\$K,V$};
 \draw[idi arrow] (memory)--(score); \draw[idi arrow] (memory)--(out);
\end{tikzpicture}
\caption{注意力把“压缩全部历史”改为“按查询重新读取历史”，代价是保存和计算位置间的相互作用。}
\label{fig:sequence-attention-query}
\end{figure}
递归模型把全部历史压入有限维状态。当预测需要重新访问某次冲击或某段工况时，这种压缩可能成为瓶颈。注意力改为按当前查询读取历史表示，其缩放点积形式为\cite{vaswani2017attention}
\begin{equation}
 \operatorname{Attention}(Q,K,V)=\softmax\left(\frac{QK^{\mathsf T}}{\sqrt{d_k}}\right)V.
\end{equation}
直接访问缩短了远距离位置之间的信息路径，但标准全局注意力的计算量随序列长度呈二次增长。Transformer 将多头注意力、位置编码、前馈网络、残差连接和归一化组合起来。因此，从 GRU 转向 Transformer 是以保存和访问历史的成本换取更灵活的信息选择，而非无条件提高预测能力。

\paragraph{注意力的逐位置解释。}
对于第$i$个查询，注意力权重为
\begin{equation}
 a_{ij}=\frac{\exp(q_i^{\mathsf T}k_j/\sqrt{d_k})}{\sum_{r=1}^{T}\exp(q_i^{\mathsf T}k_r/\sqrt{d_k})},
 \qquad o_i=\sum_j a_{ij}v_j.
\end{equation}
因此每个输出是值向量的加权平均。查询表示“我现在需要什么”，键表示“我具有什么索引”，值表示“被取出的内容”。多头注意力使用多组投影，在不同的特征子空间分别读取历史。各头可以形成不同的读取方式，但训练不保证每个头都具有可命名的分工。

\paragraph{一次可以手算的读取。}
把一个采样点或一个时间片转换成向量后，该向量也称为一个token，即模型处理的一个序列单元。对当前查询，假设两个可见历史位置的缩放得分为$0,\log3$，则softmax先取指数得到$1,3$，再除以总和得到$1/4,3/4$。若对应的标量值为$10,20$，输出就是$17.5$，而不是直接复制得分最大的$20$。键用于算权重，值用于形成输出，所以两者可以有不同维度。若第二个位置实际来自未来，必须在归一化前屏蔽它，此时权重成为$1,0$，输出为$10$。

注意力输出仍是表示，并不直接等于预测。每个头完成一次读取，多头结果按特征维拼接并投影回模型宽度；残差把原表示加回来，前馈网络再对每个位置应用同一组非线性变换。归一化控制特征尺度。这些步骤使模型既能交换位置间信息，也能加工当前位置的特征；头的具体分工由训练形成，不保证某个头必然学到某种语义。

\paragraph{缩放、归一化与访问成本的推导。}
采用逐行存放位置表示的约定。各矩阵形状为
\begin{equation}
 X\in\R^{T\times d},\quad W_Q,W_K\in\R^{d\times d_k},\quad W_V\in\R^{d\times d_v},
\end{equation}
因此$Q,K\in\R^{T\times d_k}$、$V\in\R^{T\times d_v}$。为理解缩放，暂设$q_i,k_j$的坐标彼此独立、均值为零、方差为一，则
\begin{equation}
 \E(q_i^{\mathsf T}k_j)=0,\qquad
 \Var(q_i^{\mathsf T}k_j)=\sum_{r=1}^{d_k}\Var(q_{ir}k_{jr})=d_k.
\end{equation}
因此除以$\sqrt{d_k}$使得分方差约为一，避免维度增加时softmax过早饱和。这是初始化尺度的解释，不是对训练后相关特征的严格分布假设。记缩放得分为$s_{ij}$，由商法则得到
\begin{equation}
 \frac{\partial a_{ij}}{\partial s_{ir}}=a_{ij}(\mathbf1_{j=r}-a_{ir}).
\end{equation}
当一个权重已接近一时，其余权重及大量导数接近零，模型便难以调整访问对象。softmax使$a_{ij}\geq0$且$\sum_j a_{ij}=1$，单头输出因而位于值向量的凸包内；但后续投影、残差和前馈变换不受这个凸包约束。

计算$QK^{\mathsf T}$需$O(T^2d_k)$次标量运算，计算$AV$需$O(T^2d_v)$；显式保存分数需$O(T^2)$空间。若无位置编码、无遮罩，令$P$为位置置换矩阵，则$Q'=PQ,K'=PK,V'=PV$，有$A'=PAP^{\mathsf T}$及$O'=PO$。这叫置换等变，而不是输出不随置换变化；网络本身无法从内容之外识别先后顺序。因果遮罩和位置编码正是为打破这种不足而加入的结构。

\paragraph{位置编码与因果遮罩。}
没有位置编码和遮罩的自注意力具有前述置换等变性：调换输入行会相应调换输出行，而不是让输出数表完全不变。要表达明确的先后或时间距离，需加入位置信息。正弦位置编码使用不同频率的正余弦函数；可学习位置嵌入直接为每个位置分配向量。因果遮罩令$j>i$的注意力得分为$-\infty$，从而保证第$i$个位置只能使用$1$到$i$的信息。解码时的 Key--Value (KV) cache 保存已经计算的键和值，使每次生成新 token 不必重复计算全部历史。

\subsection{Transformer 编码块怎样接到预测头}
设窗口投影与位置编码后的表示为$Z\in\R^{T\times d}$。一个前置归一化编码块可写为
\begin{align}
 U&=Z+\operatorname{MHA}(\operatorname{LN}(Z)),\\
 Z'&=U+\operatorname{FFN}(\operatorname{LN}(U)),\\
 \operatorname{FFN}(u)&=W_2\rho(W_1u+b_1)+b_2.
\end{align}
$\operatorname{LN}$是层归一化，$\operatorname{MHA}$是多头注意力（Multi-Head Attention），$\operatorname{FFN}$是逐位置前馈网络（Feed-Forward Network）。最后一式按列向量写，$W_1\in\R^{d_f\times d}$、$W_2\in\R^{d\times d_f}$，$d_f$为中间宽度。多头先各自读取历史，再拼接并投影回$d$维，以便残差相加。原始 Transformer 使用不同的归一化顺序，这里采用便于说明的前置归一化形式。

预测未来两分钟温度时，可以取最后一个位置的$d$维表示，接$2\times d$的线性投影，得到两个提前量的输出。也可以汇集整个窗口表示，再接预测头。窗口编码时所有输入都已在$t$取得，窗口内部的双向注意力不会读取$t$以后的观测。如果同时训练各位置的在线预测，则每个位置须使用因果遮罩，否则早期位置会读到它预测时尚未取得的内容。

\paragraph{多头注意力的张量形状。}
设批大小为$B$、序列长度为$T$、模型维度为$d$、头数为$H$，每头维度为$d_h=d/H$，此处假设$d$可被$H$整除；本段$H$表示头数，不是预测提前量。输入形状为$B\times T\times d$，投影后重排为$B\times H\times T\times d_h$，注意力分数形状为$B\times H\times T\times T$。最后一个$T\times T$矩阵解释了长序列的显存瓶颈：即使参数量不大，二次注意力也可能无法处理高频传感器流。

\begin{figure}[htbp]
\centering
\begin{tikzpicture}[>=Latex,node distance=0.7cm,every node/.style={font=\small}]
 \node[idi input,minimum width=1.4cm] (x) {$Z$};
 \node[idi process,right=of x,minimum width=2.5cm] (att) {归一化、注意力};
 \node[idi state,right=of att,minimum width=1.2cm] (u) {$+\ \to U$};
 \node[idi process,below=1.1cm of att,minimum width=2.5cm] (ffn) {归一化、前馈};
 \node[idi output,left=of ffn,minimum width=1.4cm] (y) {$+\ \to Z'$};
 \draw[idi arrow] (x)--(att);\draw[idi arrow] (att)--(u);
 \draw[idi arrow] (u.south)|-(ffn.east);\draw[idi arrow] (ffn)--(y);
 \draw[idi arrow,dashed] (x.north) -- ++(0,0.65) -| (u.north);
 \draw[idi arrow,dashed] (u.south) -- ++(0,-2.05) -| (y.south);
\end{tikzpicture}
\caption{前置归一化 Transformer 编码块。实线表示变换路径，虚线表示残差相加路径；注意力交换位置间的信息，前馈网络逐位置加工特征。}
\label{fig:sequence-transformer-flow}
\end{figure}

图\ref{fig:sequence-transformer-flow}沿投影、注意力和前馈变换展示编码路径，虚线则标出绕过注意力子层的残差通路；二者共同形成各位置的输出表示。
\paragraph{位置编码的选择。}
绝对位置编码把位置直接加到 token 表示中，简单但长度外推能力有限；相对位置编码直接描述 token 间距离，更适合长度变化；旋转位置编码通过复数相位或二维旋转改变查询和键，使点积携带相对位置信息。选择位置编码时要同时考虑训练长度、部署长度和缓存机制，不能只比较短序列上的损失。

\subsection{非平稳性与位置表示}
趋势、季节性、工况切换和传感器老化会使$p(x_{t+1}\mid x_{\le t})$随时间变化。差分、归一化、分解和时间特征可以减少非平稳性，也可能删除故障前的有用趋势。滚动验证比随机交叉验证更适合检验未来性能。模型更新还应区分概念漂移与测量系统改变。

\subsection{PatchTST：把局部时间片作为读取单位}
逐采样点注意力的成本随窗口长度平方增长。PatchTST（Patch Time Series Transformer）把每个传感器通道的连续片段投影成一个序列单元，再对片段做注意力。例如，一百二十个采样点按长度十、步幅十切分，得到十二个片段。单通道片段$p_j\in\R^{10}$通过$e_j=W_pp_j+b_p$映射为$d$维，其中$W_p\in\R^{d\times10}$，再加入位置信息。十二个单元的注意力分数矩阵为$12\times12$，小于逐点处理的$120\times120$。

原始 PatchTST 的一个重要设计是通道独立：各通道分别形成片段序列，并共享编码器和预测头参数。它并不通过这种注意力直接交换不同传感器的信息。预测单个温度通道时，这种共享可以利用不同通道的训练样本规律，但若预测必须依赖转速与负载的联动，还需另行设计跨通道融合或选用联合编码模型。片段投影保留多少局部细节取决于宽度和训练，片段越长也可能越难保留短冲击。规则分片不自动解决不规则采样，仍须表示实际时间间隔与缺失位置。

\section{状态空间与 Mamba：选择性地保留长历史}
\subsection{状态转移不自动等于概率估计}
状态空间模型（State Space Model，SSM）将状态演化和观测生成分开。先用线性高斯模型说明“状态表示”与“状态不确定性”的区别\cite{kalman1960filter}：
\begin{align}
 h_t&=Ah_{t-1}+Bu_t+\eta_t,\qquad \eta_t\sim\mathcal N(0,Q),\\
 x_t&=Ch_t+\varepsilon_t,\qquad \varepsilon_t\sim\mathcal N(0,R).
\end{align}
这里$u_t$为已知控制量，$Q$与$R$分别为状态噪声和观测噪声协方差。卡尔曼滤波（Kalman Filter）在状态均值之外传播协方差。确定性的 RNN 或 Mamba 层只计算一组状态数，不自动获得这种概率解释。

滤波用截至$t$的观测估计$h_t$，平滑允许利用$t$以后的观测修正过去状态，预测则推算尚未到达的状态。在线任务不能读取平滑所用的未来信息。以下保留一个滤波推导作为概率状态的对照，随后转回可训练序列层。

\paragraph{卡尔曼滤波的两步更新。}
在线线性高斯状态空间模型中，预测步骤为
\begin{align}
 \hat h_{t\mid t-1}&=A\hat h_{t-1\mid t-1}+Bu_t,\\
 P_{t\mid t-1}&=AP_{t-1\mid t-1}A^{\mathsf T}+Q.
\end{align}
观测到$x_t$后，计算创新$r_t=x_t-C\hat h_{t\mid t-1}$（新观测减去预测观测的残差，并非技术创新）和增益
\begin{equation}
 K_t=P_{t\mid t-1}C^{\mathsf T}(CP_{t\mid t-1}C^{\mathsf T}+R)^{-1},
\end{equation}
再更新$\hat h_{t\mid t}=\hat h_{t\mid t-1}+K_tr_t$。创新序列可以用于检测模型失配：若其长期偏离零或自相关显著，说明转移、噪声或传感器模型需要重新估计。

\paragraph{把滤波更新算成一个温度例子。}
设当前温度的预测均值为$20$，预测方差$P^-=4$，传感器噪声方差$R=1$，且$C=1$。读到$22$后，创新是$22-20=2$，增益为$4/(4+1)=0.8$，修正均值为$21.6$，修正方差为$(1-0.8)4=0.8$。若传感器更不可靠，$R=16$，增益降为$0.2$，均值只移到$20.4$。方差的单位是温度单位的平方；它衡量估计不确定性，不是另一条温度曲线。滤波中的“更新”是在固定系统参数下修正状态分布，与神经网络训练时更新权重不是同一操作。

\paragraph{由高斯条件化推导滤波增益。}
设$h_t\in\R^m,x_t\in\R^p,u_t\in\R^q$，则$A\in\R^{m\times m},B\in\R^{m\times q},C\in\R^{p\times m}$；$Q,P\in\R^{m\times m}$，$R\in\R^{p\times p}$。假设转移噪声与观测噪声均为零均值、彼此独立、跨时间独立，并与先前状态独立。记预测均值和协方差为$m^-,P^-$。线性变换的均值与协方差规则先给出前述预测步骤，再给出联合分布
\begin{equation}
 \begin{bmatrix}h_t\\x_t\end{bmatrix}\Bigm|x_{<t}
 \sim\mathcal N\left(
 \begin{bmatrix}m^-\\Cm^-\end{bmatrix},
 \begin{bmatrix}P^-&P^-C^{\mathsf T}\\CP^-&CP^-C^{\mathsf T}+R\end{bmatrix}\right).
\end{equation}
为看清条件化的来源，令$e=h_t-m^-$、$r=x_t-Cm^-$、$S=CP^-C^{\mathsf T}+R$，并写出修正后的误差$e^+=e-Kr$。要求它与已经利用的创新不相关，有
\begin{equation}
 \operatorname{Cov}(e-Kr,r)=P^-C^{\mathsf T}-KS=0,
 \qquad K=P^-C^{\mathsf T}S^{-1}.
\end{equation}
联合高斯使不相关等价于独立，因此$\E(h_t\mid x_t,x_{<t})=m^-+Kr$，且后验协方差为
\begin{align}
 P^+&=P^--P^-C^{\mathsf T}S^{-1}CP^-\\
 &=(I-KC)P^-(I-KC)^{\mathsf T}+KRK^{\mathsf T}.\label{eq:sequence-joseph-covariance}
\end{align}
式\eqref{eq:sequence-joseph-covariance}称为Joseph形式，更适合在有限精度计算中保持对称性和半正定性。这里要求$S$可逆，实际实现应解线性方程而非显式求逆。若$p=m=1,C=1$，则$K=P^-/(P^-+R)$：观测越噪，修正越小；状态预测越不确定，越相信新观测。在线性高斯且参数正确时，$r^{\mathsf T}S^{-1}r$服从自由度为$p$的卡方分布；用它设阈值仍需检查噪声、参数估计与重复告警假设，不能直接把名义尾概率当成事件级误报率。

\subsection{从固定状态空间到 Mamba 的选择性记忆}
\label{sec:sequence-mamba}
前面的状态空间模型（State Space Model，SSM）用概率分布描述状态与观测，并通过卡尔曼滤波估计不确定性。本节沿用状态转移的语言，但把 SSM 用作可训练的确定性序列层：它输出表示而非状态后验，不自动携带协方差或预测区间。先说明固定线性系统为何能形成长记忆，再解释 Mamba 如何根据输入选择写入、保留和读取的信息。以下讨论原始 Mamba 的选择性 SSM 与模块结构，不把后续 Mamba-2 的算法细节混入其中。

\paragraph{连续动力学怎样成为离散记忆。}
以$\tau$表示连续时间，$t$仍表示离散位置。设状态$h(\tau)\in\R^m$、输入$u(\tau)\in\R^q$、层输出$v(\tau)\in\R^r$，连续系统为
\begin{equation}
 \frac{d h(\tau)}{d\tau}=A_c h(\tau)+B_c u(\tau),\qquad
 v(\tau)=Ch(\tau)+Du(\tau).
 \label{eq:sequence-ssm-continuous}
\end{equation}
其中$A_c\in\R^{m\times m}$、$B_c\in\R^{m\times q}$、$C\in\R^{r\times m}$、$D\in\R^{r\times q}$均不随位置变化，$D$为输入到输出的直接通路。这里$u_t$是由观测$x_t$提取的输入特征，不限于前文滤波模型中的控制量；$v_t$是中间表示，不与任务标签$y_t$混用。式\eqref{eq:sequence-ssm-continuous}是线性时不变（Linear Time-Invariant，LTI）系统。

在长度为$\Delta>0$的每个区间内保持输入不变，即零阶保持，积分因子法给出
\begin{align}
 h_t&=\bar A h_{t-1}+\bar B u_t,\qquad v_t=Ch_t+Du_t,
 \label{eq:sequence-ssm-discrete}\\
 \bar A&=\exp(\Delta A_c),\qquad
 \bar B=\int_0^\Delta\exp(sA_c)B_c\,ds.
 \label{eq:sequence-ssm-zoh}
\end{align}
只有$A_c$可逆时，才可把式\eqref{eq:sequence-ssm-zoh}的积分写为$A_c^{-1}(\bar A-I)B_c$；积分形式始终有效。反复代入式\eqref{eq:sequence-ssm-discrete}可得
\begin{equation}
 v_t=C\bar A^t h_0+\sum_{j=1}^t C\bar A^{t-j}\bar B u_j+Du_t.
 \label{eq:sequence-ssm-convolution}
\end{equation}
式\eqref{eq:sequence-ssm-convolution}表明，固定线性系统既可递推，也可使用只依赖时间差的卷积核$G_k=C\bar A^k\bar B$计算。S4通过结构化参数高效表示和计算长核；卷积实现的复杂度取决于具体算法，不能将其一概说成严格线性。若$A_c$的特征值实部为负，则$\bar A$的特征值模小于一，长期记忆衰减，但非正规矩阵仍可能短期放大扰动。固定核对同一时间差使用相同权重，不能仅因某次冲击内容重要就改变记忆规则，这正是选择性更新要解决的限制。

\paragraph{输入依赖的步长、写入与读取。}
为明确维度，先看一个标量输入通道$u_t\in\R$，其状态$h_t\in\R^m$。令$\xi_t\in\R^d$为当前位置经过因果局部处理后的特征，定义
\begin{align}
 \Delta_t&=\operatorname{softplus}(w_\Delta^{\mathsf T}\xi_t+b_\Delta)>0,
 \label{eq:sequence-mamba-selection}\\
 B_t&=W_B\xi_t\in\R^m,\qquad C_t=W_C\xi_t\in\R^m,\qquad
 A_c=\diag(-\exp(a_1),\ldots,-\exp(a_m)).\nonumber
\end{align}
这里$w_\Delta\in\R^d$、$W_B,W_C\in\R^{m\times d}$以及$a_i$是训练得到、推理时共享的参数；$\operatorname{softplus}(s)=\log(1+e^s)$。式\eqref{eq:sequence-mamba-selection}是便于推导的单通道参数化，实际模块可用低秩投影产生各通道步长，并在通道间共享部分$B_t,C_t$。输入依赖的是$\Delta_t,B_t,C_t$，不是所有参数；尤其$A_c$和直接通路系数$D$不必随输入变化。

逐步冻结这些系数，精确零阶保持对应
\begin{align}
 \bar A_t&=\exp(\Delta_t A_c),\qquad
 \bar B_t^{\mathrm{ZOH}}=\int_0^{\Delta_t}\exp(sA_c)B_t\,ds,
 \label{eq:sequence-mamba-zoh}\\
 h_t&=\bar A_t h_{t-1}+\bar B_t u_t,\qquad
 v_t=C_t^{\mathsf T}h_t+Du_t.
 \label{eq:sequence-mamba-recurrence}
\end{align}
式\eqref{eq:sequence-mamba-recurrence}中的$\bar B_t$必须与所选离散化一致。原始 Mamba 的官方选择性扫描实现对输入项采用$\bar B_t=\Delta_t B_t$，同时保留$\bar A_t=\exp(\Delta_t A_c)$；它并非直接使用式\eqref{eq:sequence-mamba-zoh}的精确输入积分。两者在$\|\Delta_t A_c\|$较小时一阶相符，不能将简化输入项写成精确零阶保持的恒等式。

对第$i$个状态分量，保留系数为$\exp(-\Delta_t\exp(a_i))$。较小的$\Delta_t$使旧状态较少衰减，较大的$\Delta_t$使更新更快；$B_t$控制写入方向和强度，$C_t$控制当前读出。步长因此是一种学习到的记忆时间尺度，不天然等于传感器的物理采样间隔。对于不规则采样，还需显式提供实际时间间隔或约束连续时间模型。选择也不是“自动丢弃噪声”的保证：是否保留短时冲击取决于训练目标与数据，必须通过事件级验证检查。

尽管式\eqref{eq:sequence-mamba-recurrence}对给定系数下的旧状态仍是线性的，系数本身依赖输入，整个输入到输出映射已不是 LTI 系统。历史输入的权重同时依赖其后的选择，不能再用式\eqref{eq:sequence-ssm-convolution}中那个固定的时间差卷积核替代。

\paragraph{先手算两步，再理解扫描。}
取单维$h_0=0$，第一步$F_1=0.5,b_1=2$，第二步$F_2=0.8,b_2=1$。依次递推得到$h_1=2,h_2=0.8\times2+1=2.6$。合并成一步则是$h_2=(0.8\times0.5)h_0+(0.8\times2+1)=0.4h_0+2.6$。第三步既可先与第二步合并，也可后接到前两步的结果上，答案相同；交换前两步却得到$0.5\times1+2=2.5$，说明顺序不能交换。这里的$b_t$是已经计算好的写入量，不是神经网络所有层都共用的固定偏置。

离散化也可以手算：取$A_c=-1,B_t=1,u_t=2,h_{t-1}=0,\Delta_t=\log2$，则保留系数为$1/2$。精确积分的输入系数为$1-e^{-\Delta_t}=1/2$，得到$h_t=1$；简化输入系数为$\Delta_t$，得到$h_t=2\log2\approx1.386$。这不是计算错误，而是两种离散化在步长不小时确实不同。学习参数可以适应所用实现，但不能因此把两种公式称为相等。

\paragraph{为何能递推推理，也能并行训练。}
记$F_t=\bar A_t$、$b_t=\bar B_t u_t$，每步状态更新都是仿射映射$h\mapsto F_t h+b_t$。定义按时间先后合并两个映射的运算
\begin{equation}
 (F_2,b_2)\star(F_1,b_1)
   =(F_2F_1,\ F_2b_1+b_2).
 \label{eq:sequence-mamba-scan}
\end{equation}
函数复合的结合性使式\eqref{eq:sequence-mamba-scan}满足结合律，单位元为$(I,0)$；它通常不满足交换律，不能打乱时间顺序。对前缀$(F_t,b_t)\star\cdots\star(F_1,b_1)$做并行扫描，即可得到所有$h_t$。由于选择参数由输入特征而非尚未求得的$h_{t-1}$生成，训练时可先并行计算各位置系数，再合并前缀；这与一般非线性 RNN 的顺序状态依赖不同。理想并行前缀算法具有$O(\log T)$深度，但总工作量仍随$T$增长；实际硬件上的分块、访存和同步决定实际延迟，不能据此宣称运行时间只随$\log T$增长。

\paragraph{从选择性 SSM 到 Mamba 模块。}
原始 Mamba 不只是一个状态递推式。输入先经线性投影分成内容分支与门控分支；内容分支经过短核因果深度卷积和 SiLU 激活，再进入选择性 SSM。这里深度卷积指每个通道分别沿时间使用短核，SiLU为$\operatorname{SiLU}(a)=a\sigma(a)$，两者分别负责局部混合与非线性变换。扫描结果与另一分支的 SiLU 输出逐元素相乘，经输出投影后接入残差路径；堆叠模块配合归一化构成深层网络。因果卷积只读取当前及过去的位置，承担局部混合，SSM承担长历史压缩，通道投影混合不同传感器特征。流式部署必须同时维护 SSM 状态和短卷积缓存。输出还须接任务头，才能成为未来轨迹、故障概率或分位数预测；选择性状态本身不是概率置信度。

\paragraph{与 Transformer 的计算和记忆取舍。}
设层内通道数为$d$，每通道状态维度为$m$。对角转移下，选择性扫描核心的总工作量为$O(Tdm)$，流式状态占用为$O(dm)$，与已处理历史长度无关；固定短卷积核还需要与核宽有关的缓存。密集通道投影通常另需$O(Td^2)$工作，不能把整个模块的成本仅写成$O(T)$而省略宽度与状态大小。训练仍需存储或重算随序列长度增长的激活，硬件感知扫描通过融合、分块与重计算减少访存，不意味着训练显存为常数。

相较之下，标准全局注意力的交互计算为$O(T^2d)$，自回归 KV 缓存随历史增长到$O(Td)$，每个新位置还需读取历史键和值。FlashAttention 等精确算法可避免显式存储完整注意力矩阵，但不消除全局交互的二次算术量。Mamba 在长流与受限缓存下有结构性优势，Transformer 则能按当前查询重新访问具体历史位置；压缩进有限状态的证据一旦丢失，Mamba 无法像注意力那样直接回读原记录。精确事件检索、短序列、小批量或缺少优化扫描内核时，不能预设 Mamba 更准确或更快；混合注意力与 SSM 也是可比较的选择。

\paragraph{工业实例：保留冲击还是保留工况。}
沿用轴承案例，可将因果振动片段特征、当前温度、转速、负载、缺失掩码与实际时间间隔组成$x_t$，经投影和局部卷积形成$\xi_t$。稳定工况中，一部分状态维度可学习较慢衰减以积累磨损趋势；负载切换时，选择机制可调整写入，避免把正常升温直接解释为退化；稀疏冲击是否进入长期记忆则由故障目标约束。这只是机制解释，不是已经观察到的实验结果，更不能把某个门值当作物理故障原因。

在相同设备划分、时间边界、预测提前量和参数预算下，应比较 GRU、Transformer 与 Mamba，并消融输入依赖的$\Delta_t,B_t,C_t$，检查收益究竟来自选择机制还是局部卷积与额外容量。除轨迹误差外，报告短冲击召回、事件级误报、流式延迟和峰值显存。每台设备独立维护状态；跨设备拼接、长缺失后的恢复和模型更新必须按预定规则重置或预热，训练时也要模拟这些边界。这样，式\eqref{eq:sequence-mamba-recurrence}中的长记忆才与上线时真正可用的历史一致。

\section{把历史表示接到预测与异常任务}
\subsection{直接多步预测与递归预测}
前面的模型都先形成表示，输出头才决定预测什么。直接多步预测用$\hat y_{t+1:t+H}=Wr_t+b$一次生成$H$个未来值，其中$W\in\R^{H\times d}$、$r_t\in\R^d$。可用
\begin{equation}
 \mathcal L_{\mathrm{point}}=\frac{1}{NH}\sum_{n=1}^N\sum_{h=1}^H(\hat y_{t_n+h}-y_{t_n+h})^2
\end{equation}
训练，$N$为训练窗口数，$t_n$为第$n$个窗口终点。各提前量的输出共享表示，但有不同的投影行。平方损失对应条件均值估计，不能单凭点预测给出区间。

递归预测每次只预测一步，再把预测值反馈给下一步。以标量温度为例，$\hat y_{t+1}=F_\theta(y_{t-T+1:t})$，第二步将窗口移为$(y_{t-T+2},\ldots,y_t,\hat y_{t+1})$。若模型还需要未来外生变量，必须使用已经知道的计划值或另行预测。训练时用真实历史作为递归输入称为教师强制（Teacher Forcing），部署则会遇到自己的预测误差。因此必须用实际部署的反馈方式评价多步结果。

为说明误差如何传播，暂设真实状态$x_{t+1}=f_*(x_t)+\eta_{t+1}$，递归预测$\hat x_{t+h}=f_\theta(\hat x_{t+h-1})$且$\hat x_t=x_t$。在所考察区域假设$\|f_\theta(a)-f_\theta(b)\|\leq L\|a-b\|$、$\|f_\theta(x)-f_*(x)\|\leq\delta$、$\|\eta_t\|\leq\epsilon$。插入并减去$f_\theta(x_{t+h-1})$，得到
\begin{align}
 e_h&=\|\hat x_{t+h}-x_{t+h}\|\leq Le_{h-1}+\delta+\epsilon,\\
 e_h&\leq(\delta+\epsilon)\sum_{j=0}^{h-1}L^j.
\end{align}
$L=1$时上界随$h$线性增长，$L>1$时可能更快放大，$L<1$时此上界有界。这是有界噪声条件下的推导，高斯噪声不能直接套用统一确定性$\epsilon$。直接输出避免反馈自身误差，但仍受表示能力和未来不可预测性限制。

\subsection{概率输出怎样表达未来的不确定性}
若每个提前量输出$0.1,0.5,0.9$三个分位数，未来两分钟的输出形状为$2\times3$。其中中位数描述条件分布的中间位置，上下分位数给出名义覆盖率为$80\%$的区间。分位数损失（Pinball Loss）为
\begin{equation}
 \ell_\tau(y,a)=\max\{\tau(y-a),(\tau-1)(y-a)\},\qquad 0<\tau<1.
\end{equation}
$y$为实际未来值，$a$为对应分位数预测。取$\tau=0.9,y=10$，预测$8$的损失为$1.8$，预测$12$为$0.2$，因此高分位数更重罚低估。

固定历史，令未来标量的条件分布函数为$F$。在期望损失有限且$F$连续时，
\begin{align}
 R(a)&=(1-\tau)\int_{-\infty}^a(a-y)\,dF(y)+\tau\int_a^\infty(y-a)\,dF(y),\\
 R'(a)&=(1-\tau)F(a)-\tau(1-F(a))=F(a)-\tau.
\end{align}
最优点满足$F(a)=\tau$，有原子质量时以次梯度条件$F(a^-)\leq\tau\leq F(a)$代替。训练后须逐提前量检查覆盖率和区间宽度，过宽的区间虽然容易覆盖，却可能没有使用价值。独立分位数可能交叉，模型失配与漂移也会改变覆盖率。各时刻分别覆盖不等于整条轨迹同时覆盖。

未来故障概率使用另一种输出。令$c_t\in\{0,1\}$表示$(t,t+H]$内是否发生指定故障，输出$p_t=\sigma(w^{\mathsf T}r_t+b)$，可用二元交叉熵$-c_t\log p_t-(1-c_t)\log(1-p_t)$训练。它的标签是未来事件，不能用当前异常分数直接替代。概率输出仍需检查校准，少数故障事件不能仅用准确率评价。

\subsection{异常评分：把新观测与正常模式比较}
预测残差检测可预测模式的偏离。若预测头还输出正常数据的条件分布，则可用负对数似然评分。例如单通道正态预测$\mathcal N(\mu_t,\sigma_t^2)$在读到$y_t$后给出
\begin{equation}
 s_t^{\mathrm{NLL}}=\frac{(y_t-\mu_t)^2}{2\sigma_t^2}+\log\sigma_t+\tfrac12\log(2\pi),\qquad\sigma_t>0.
\end{equation}
$\mu_t,\sigma_t$须在读到$y_t$之前生成，传感器单位保持固定。相同误差在正常波动小的工况下得分更高。方差估计失真也会改变异常排序，因此需要独立的正常和异常验证数据。

还可以训练窗口重构模型，把编码器表示经解码器还原为$\tilde X_t$，以重构误差衡量偏离。它可以使用窗口中的当前观测，判断对象是已经到达的窗口。模型也可能重构异常样本，尤其是训练集中混入异常或容量过大时，低重构误差不能保证正常。本章以预测残差为主要算例，所有异常分数都需要验证其与实际事件的关系。

\subsection{时间边界与工业验证}
\paragraph{滚动预测的完整设置。}
设传感器以一分钟为间隔采样，使用过去$T=120$步预测未来$H=30$步。样本构造应令输入窗口为$[t-T+1,t]$，标签窗口为$[t+1,t+H]$，并保证同一故障事件不会同时出现在训练和测试窗口。若窗口高度重叠，随机切分会使相邻片段泄漏到不同集合；更稳妥的做法是按时间块或设备划分，并在块之间留出间隔。

评价时分别计算$h=1,5,15,30$的误差，观察模型是否只在短提前量有效。若采用递归预测，必须在每一步输入模型自己的预测；若采用直接多输出头，则应说明每个提前量是否共享参数。除 MAE 外，还应报告故障前提前量、告警持续时间和每台设备每天误报数，因为这些量直接决定维护团队的工作量。

\paragraph{从预测到异常告警。}
轨迹预测误差只有在真实观测到达后才能计算，不能用未来误差触发当前预警。设$\hat x_{t\mid t-h}$及$\hat\sigma_{t\mid t-h}$是在$t-h$时刻对$t$的预测和标准差。对一个传感器通道，时刻$t$可计算的标准化残差为
\begin{equation}
 s_t=\max_{1\leq h\leq H}\frac{|x_t-\hat x_{t\mid t-h}|}{\hat\sigma_{t\mid t-h}+\epsilon}.
\end{equation}
其中$\epsilon>0$须与标准差使用相同单位。即使单个提前量的标准化残差近似正态，取绝对值再对多个提前量取最大值后，也不能直接套用标准正态阈值；应按实际组合规则在验证集上定阈值。这个分数检测当前观测相对于历史预测的偏离，不等于未来故障概率。若要提前预警，还需验证偏离与后续故障的关系，或直接训练未来事件概率输出头。
为了避免一个瞬时噪声触发告警，可以要求连续$m$个窗口超过阈值，或在事件级别合并相邻告警。阈值和$m$只能在验证数据上根据漏报成本、误报预算和最小提前量确定。测试集只用于一次性报告，不能在看到故障时间后反复调整规则。

传感器掉线时，要先区分设备真的异常，还是根本没有可靠读数。直接填零可能在曲线上制造一次突变。可以用缺失掩码注明哪些值没测到；若沿用上次读数，还应限制最多沿用多久。长时间缺失应触发数据质量告警，并按预定规则暂停判断或转人工处理。这样才不会把通信故障误报成设备故障。

\section{回到设备监测：怎样验证历史表示有用}
\paragraph{序列模型的基线与消融。}
持久性预测$\hat x_{t+h}=x_t$是许多连续传感器任务的强基线；具有周期性的过程还应加入季节性复制和线性外推。若深度模型不能稳定超过这些基线，应先检查窗口构造、目标定义和数据泄漏，而不是立即增加层数。消融实验可以依次去掉时间特征、缺失掩码、注意力、残差或多步损失，以判断性能来自真正的结构还是额外的参数量。

\paragraph{不规则采样与时间间隔。}
许多工业事件并非严格等间隔。直接把事件编号当作时间步，会把十秒间隔和十分钟间隔视为同样的动态变化。可以显式输入$\Delta t$，或使用连续时间状态转移；但填补缺失时间点也会改变噪声结构。实验应分别比较规则重采样、时间间隔特征和事件驱动建模，并报告重采样造成的延迟和短时故障损失。

\paragraph{多变量序列的因果边界。}
多个传感器之间可能存在同步采集、控制回路和滞后影响。相关性注意力权重不能直接解释为因果关系，因为共同工况变量会同时改变多个通道。若预测目标是控制动作，还应区分可用的外生变量、动作之后才产生的变量和部署时不可获得的维修记录。特征可用性应按时间戳审计，并在数据字典中固定定义。

\paragraph{概率预测的评价。}
分位数预测使用 pinball loss，完整分布预测还可以使用连续排序概率得分。评价时既要检查尖锐性，也要检查校准：过宽的预测分布可能轻易覆盖真实值，却没有调度价值。对多步预测，应分别评估每个提前量，并观察区间是否随着预测跨度合理变宽。若预测分布出现不合理的收缩，通常意味着模型忽略了过程噪声或训练目标过度追求均值误差。

以下材料分别对应递归状态、注意力和选择性状态空间。阅读时应回到本章的共同问题：模型保留了哪些历史，输出服务于哪个任务，评价时使用的信息是否与上线一致？
\section*{辅助学习材料}
\begingroup
\emergencystretch=3em
\begin{itemize}
\item \textbf{CS224N: Natural Language Processing with Deep Learning}；作者/平台：Stanford University；英文；课程讲义、作业与视频。重点看 self-attention、Transformers 和 evaluation，帮助理解本章注意力、位置编码与评价边界；课程内容随学年更新。\href{https://web.stanford.edu/class/cs224n/}{课程主页}。
\item \textbf{NLP From Scratch: Classifying Names with a Character-Level RNN}；作者/平台：Sean Robertson / PyTorch；英文；代码教程。阅读字符编码、循环网络、训练与混淆矩阵部分，观察隐藏状态如何逐字符更新以及如何完成序列分类。\href{https://docs.pytorch.org/tutorials/intermediate/char\_rnn\_classification\_tutorial.html}{在线阅读}。
\item \textbf{The Illustrated Transformer}；作者/平台：Jay Alammar / 作者博客；英文；图解教程。对照本章查询、键、值、多头注意力和位置编码，沿编码器与解码器的信息流理解 Transformer；注意区分解码器因果遮罩与编码器的双向访问。\href{https://jalammar.github.io/illustrated-transformer/}{在线阅读}。
\item \textbf{A Visual Guide to Mamba and State Space Models}；作者/平台：Maarten Grootendorst / 作者博客；英文；图解教程。配合第\ref{sec:sequence-mamba}节，从连续与离散 SSM、递推与卷积读到输入依赖选择和扫描；并行扫描依赖结合律而非交换律，时间顺序不能改变。\href{https://www.maartengrootendorst.com/blog/mamba/}{在线阅读}。
\item \textbf{Mamba}；作者/平台：state-spaces / GitHub 官方项目文档；英文；实现与使用说明，非论文阅读条目。重点看 Usage 下的 Selective SSM、Mamba Block，沿文档链接核对扫描接口与模块实现；区分原始 Mamba 与后续版本，部署前检查平台和加速内核要求。\href{https://github.com/state-spaces/mamba\#usage}{使用文档}。
\end{itemize}
\endgroup
\paragraph{自测与参考答案。}
在$t$时刻，能否用刚生成的$t+5$预测与其真实值之差触发告警？若把训练计算图每20步切断，是否必须每20步清零状态？参考答案：未来真实值尚未到达，不能计算该残差；只能使用当前观测与过去对当前的预测比较，或使用另行训练的未来事件概率。截断只切断梯度连接，连续序列可保留状态数值；换设备或独立序列才按边界规则重置。

\paragraph{本章小结。}
序列建模先确定截至当前可取得的观测，再选择历史表示方法。RNN 逐步更新状态，LSTM 与 GRU 控制旧信息的保留和新信息的写入。TCN 通过因果空洞卷积组合固定感受野，Transformer 通过查询重新读取历史，PatchTST 将读取单位改为局部片段。Mamba 让输入控制状态更新，并用结合性的仿射复合支持并行扫描。这些机制分别限制了模型能访问的历史、计算路径和缓存需求，没有一种结构对所有任务都更合适。

表示之后，未来轨迹使用数值或分位数头，未来故障使用事件概率头，当前异常则比较已到达的观测与正常模式。概率状态滤波、确定性隐状态和概率预测头具有不同含义。多步预测要按实际反馈策略评价，异常告警要按观测到达时间计算，模型比较要固定时间边界、设备划分和资源约束。

振动曲线可以提示设备发生了变化，但通常不能指出表面裂纹在哪里。下一章转向图像，用卷积、区域匹配和像素分割寻找缺陷位置。两章都要求检查输入与标签的时间或空间边界，但序列章节的输出是未来数值、事件概率或已到达观测的异常分数，不能拿定位成绩替代这些指标。
